\begin{proposition}[Local complexes and every noncusp torsion length]\hypertarget{tag-01SR}{}\label{inv:D4.4}
\StatementRecord{inv:D4.4}{proved}{01SR}{\SeriesRef{inv:D4.1}, \SeriesRef{inv:A2.3}}

Let $z\notin\Z$. At a regular resonance $t_*$ in \eqref{pub:I:eq:D-15}, with
$A=\C\{t-t_*\}$, the two stalks are
\begin{equation}
 (R^1f_*(T_X\otimes M_z))_{t_*}=(A/(t-t_*))^2,
 \qquad (R^2f_*(T_X\otimes M_z))_{t_*}=A/((t-t_*)^2).
\label{pub:I:eq:D-18}
\end{equation}
At $p_1$, write $z=n+k(3\beta_1+4-2\rho)$, $k\ne0$, when this is
possible. Both coarse stalks are killed by $t$ and have length
$1+\epsilon_n$, with $\epsilon_n$ as in \ref{inv:D4.1}.
At $p_2$, write $z=n+k(2\beta_2+3-3i)$, $k\ne0$, when possible.
Both coarse stalks are killed by $t-1$ and have length one for
$n\equiv1,2\pmod4$ and length zero for $n\equiv0,3\pmod4$.
If the relevant lattice representation does not exist, both stalks
at that finite fibre are zero.
\end{proposition}


\begin{proof}
\begin{enumerate}[label=\arabic*.,leftmargin=*]
\item We first give the comparison that turns the local calculation into
a computation of the sheaves. On a simply connected disc of smooth
torus fibres, including a finite root cover, retain the real marking
and the holomorphic period matrix. Translation by $\Pi(s)v$, for
fixed $v\in\R^4$, is holomorphic on the total family. After
subtracting the integer cocycle of a resonance, the line bundle has
constant lattice multipliers. Multiplying the lifted real translation
by its constant factor of automorphy makes the lattice act as the
identity and gives a genuine compact real torus action on the line
bundle. The same action acts on the tangent bundle by differentiation.
It commutes with the full Dolbeault differential.

\item Averaging over this compact torus projects the relative Dolbeault
complex onto its invariant part. Choose the lift so that the central
trivial line has trivial character. On every nonzero Fourier character,
the relative Dolbeault symbol is exterior multiplication by a
nonzero covector; contraction by a dual covector is a contracting
homotopy. After shrinking the disc, these nonzero covectors have a
uniform positive lower bound for bounded Fourier frequency and grow
linearly for unbounded frequency, because the real period matrices
and their inverses stay bounded. The homotopies therefore preserve
smooth Fourier series, with all derivatives, and are uniform on
compact subsets of the disc. The tangent extension is triangular
with respect to its vertical and horizontal parts; its nilpotent
off-diagonal term is treated by the finite homological perturbation
series, so every nonzero Fourier complex is still contractible.
These homotopies exhibit averaging as a quasi-isomorphism.

\item The chain-level comparison uses the following frames. Write the normalized period matrix as $[Z(s),I]$
and use the real coordinates $\zeta=Z(s)x+y$, with
$x,y\in\R^2/\Z^2$. Set
\[
 G_j=(dx_j)^{0,1},\qquad e_i=\partial_{\zeta_i},\qquad
 h=\partial_s+\sum_i(Z'(s)x)_i e_i.
\]
These are invariant under the real translations. Since each $dx_j$
is closed on the total real product, $\bar\partial G_j=0$.
The forms $G_j$ frame the relative invariant $(0,1)$ forms because
$\operatorname{Im}Z$ is invertible, and
\[
 \bar\partial h=\sum_{i,j}Z'_{ij}(s)e_iG_j,
 \qquad (dy)^{0,1}=-Z(s)(dx)^{0,1}.
\]
For the smooth line frame defined by its flat character, the last
identity gives $\bar\partial\ell=2\pi i(a_1G_1+a_2G_2)\ell$,
with $a_1,a_2$ in \eqref{pub:I:eq:D-20}; there is no additional base term.
All displayed coefficients are holomorphic in $s$. The full
invariant complex is thus the ordinary base Dolbeault resolution
of the following finite complex of holomorphic bundles:
\[
 C^j=(V\oplus L)\otimes\bigwedge^jG,\quad0\le j\le2,
 \qquad
 d(v\otimes\eta)=v\otimes a\wedge\eta,
\]
\begin{equation}
 d(h\otimes\eta)=h\otimes a\wedge\eta+\kappa\wedge\eta.
\label{pub:I:eq:D-19}
\end{equation}
Here $V$ is the rank-two vertical bundle, $L$ the horizontal line,
$G=R^1\OO$ on the root chart, and $\kappa\in V\otimes G\otimes L^*$
is the tangent-extension cocycle. The matrix is obtained
by differentiating the period translations. In the raw period and
constant-cocycle frames it is
\begin{equation}
 \kappa=\begin{pmatrix}6\mu'&\tau'\\\beta_c'&\mu'\end{pmatrix},
 \qquad
 a=(z-n+6q\mu+r\beta_c,-p+q\tau+r\mu).
\label{pub:I:eq:D-20}
\end{equation}
Let $\lambda=2\pi i$. The Dolbeault exponential convention gives
the differential with Picard vector $\lambda a$ and tangent term
$\kappa$. Its explicit chain isomorphism to \eqref{pub:I:eq:D-19} is, in degree $j$,
\[
 F^j|_{V\otimes\bigwedge^jG}=\lambda^{-j}\operatorname{id},
 \qquad
 F^j|_{L\otimes\bigwedge^jG}=\lambda^{-j-1}\operatorname{id}.
\]
Direct substitution gives $F^{j+1}d_{\lambda a,\kappa}
=d_{a,\kappa}F^j$ on both summands. The inverse multiplies them
by $\lambda^j$ and $\lambda^{j+1}$, respectively. These constant
maps also commute with the finite cyclic action. The displayed formula also follows directly by
subtracting the holomorphic period rows from the constant lattice
cocycle. Together these formulas identify the chain map and its tangent extension.
For the nonzero Fourier characters, use the filtration by the number
of base $(0,1)$ factors: the relative contractions make the associated
graded complexes acyclic, and the two-step filtration then proves
acyclicity of the full complexes. Thus the contraction need not
commute with the base Dolbeault operator. The finite cyclic group
normalizes the real translation group, so averaging and this
comparison are equivariant; exact invariants compute the coarse
multiple-fibre stalks.

\item At a regular resonance, $a'=\kappa^t(q,r)^t$ and
$(q,r)\ne(0,0)$. The determinant of $\kappa$ is nonzero there by
\cref{inv:A2.1}. Therefore $a$ has a simple zero. Choose its nonzero leading
direction as the second frame of $G$ and use invertibility of
$\kappa$ to choose the two vertical frames. Holomorphic invertible
changes reduce \eqref{pub:I:eq:D-19} to
\begin{equation}
 a=(0,s),\qquad \kappa=\begin{pmatrix}0&1\\1&0\end{pmatrix}.
\label{pub:I:eq:D-21}
\end{equation}
The second cohomology is generated by $h\gamma\delta$ and is killed
by $s^2$; the relations eliminate the two vertical generators.
The first cohomology is generated by $e_1\delta,e_2\delta$, each
killed by $s$. These are equalities of modules over $\C\{s\}$
and prove \eqref{pub:I:eq:D-18}.

\item There is an additional joint coefficient to check at the order-three
point. Put $t=s^3$, $\nu=\mu-(2-\tau)/3$ and
$b=\beta_c+4\mu+(2/3)\tau-b(0)$. The order-three transformation
laws give $\nu=\nu_2s^2+O(s^3)$ and $b=b_3s^3+O(s^4)$.
The elliptic calculation in \ref{inv:A2.3} gives, with elliptic scale
$f(s)=F s^2+O(s^3)$,
\[
 \nu_2^2=9F^2/512.
\]
For clarity we evaluate the other coefficients with the original
Weierstrass normalization. Since $x=f^2\wp$,
$y=f^3\wp'/2$, one has
$f^4g_2=12t^3(t-1)$ and $f^6g_3=-8t^4(t-1)^2$.
The identities $g_2=4\pi^4E_4/3$, $g_3=8\pi^6E_6/27$ and
$E_4'(\rho)=-2\pi iE_6(\rho)/3$ give
\[
 E_6(\rho)=-\frac{27}{\pi^6F^6},\qquad
 \tau(s)=\rho+\tau_1s+O(s^2),\quad \tau_1=\frac{i\pi F^2}{2}.
\]
The period determinant identity in \cref{inv:A2.1} gives
\[
 K(s)=6\mu'^2-\tau'\beta_c'=K_2s^2+O(s^3),
 \quad K_2=\frac{3F^2}{4}.
\]
Equivalently $K=6\nu'^2-\tau'b'$, so
$K_2=24\nu_2^2-3\tau_1b_3$ and $b_3=7i/(32\pi)$.

\item For a nonintegral resonance the relation is
$(q,r,p)=\pm(k,3k,2k)$, $k\ne0$; the sign only changes the
following nonzero scalar. After a constant change of the two
cocycle frames, its Picard section is $3k(\nu,b)$, and the
transpose period derivative becomes
\[
 \begin{pmatrix}\tau'&\nu'\\6\nu'&b'\end{pmatrix}.
\]
Eliminating the unit $\tau'$ leaves the joint transverse component
$3k(b-6\nu\nu'/\tau')$. Its coefficient of $s^3$ is
\begin{equation}
 -\frac{k(K_2+12\nu_2^2)}{\tau_1}
       =\frac{123ik}{64\pi}\ne0.
\label{pub:I:eq:D-22}
\end{equation}
This nonzero joint coefficient determines the two transverse
orders in the following Smith presentation.

\item Choose the second cocycle frame to be the Picard section divided by
$s^2$. Equivariance and \eqref{pub:I:eq:D-22} give a matrix with a unit in its
lower-right entry, first-order off-diagonal terms with nonzero
coefficients, and a second-order upper-left term. Subtracting a
holomorphic multiple of the second vertical row removes the last
term; the nonzero determinant coefficient and \eqref{pub:I:eq:D-22} keep both
off-diagonal coefficients units. Rescaling the equivariant frames
therefore reduces \eqref{pub:I:eq:D-19} to
\begin{equation}
 a=(0,s^2),\qquad \kappa=\begin{pmatrix}0&s\\s&1\end{pmatrix}.
\label{pub:I:eq:D-23}
\end{equation}
Its cohomology is
\begin{equation}
 H^1=\C\{s\}/(s^2)\oplus\C\{s\}/(s^3),\qquad
 H^2=\C\{s\}/(s^2)\oplus\C\{s\}/(s^3).
\label{pub:I:eq:D-24}
\end{equation}
For example in degree one put
$u=e_1\delta$, $v=e_2\gamma+s h\delta$; the remaining cycle
$e_2\delta$ is eliminated by the boundary
$s u+e_2\delta+s v$. The two surviving relations are
$s^2u=0,s^3v=0$. In degree two eliminate $e_2\gamma\delta$;
the remaining generators $e_1\gamma\delta,h\gamma\delta$
have annihilators $s^2,s^3$ respectively.

\item We use the following character convention.
A frame has invariant exponent $w$ when its coefficient satisfies
$g(\zeta s)=\zeta^w g(s)$; its pullback-frame weight is the negative
of $w$. For the trivialized twist in \eqref{pub:I:eq:D-16}, these exponents are
$n$ at $p_1$ and $3n$ at $p_2$. This agrees with $M_n$ in \ref{inv:D4.1}
and follows by evaluating its integer cocycle on $v_1,v_2$.
At $p_1$ the two generators in $H^1$ of \eqref{pub:I:eq:D-24} have invariant
exponents $n+2,n+2$, and the two generators in $H^2$ have exponents
$n+2,n$. The length-three summand has one invariant coefficient
for every $n$; the length-two summand has one precisely when
$n\not\equiv0\pmod3$. The coarse parameter is $s^3$, so all these
invariants are killed by it. This proves the asserted order-three
stalks with the full nilpotent layers included.

\item At $p_2$ the determinant is a unit. The Picard section has a simple
zero, so \eqref{pub:I:eq:D-21} applies on the order-four cover. The invariant
exponents of the two degree-one generators are $3n+1,3n+2$;
the cyclic degree-two module of length two has generator exponent
$3n+2$. In both degrees its invariant dimension is one exactly for
$n\equiv1,2\pmod4$. The coarse parameter $s^4$ kills the modules.
This proves the stated order-four calculation. \end{enumerate}
\end{proof}